\chapter{Vectors, Bases, Matrices, and Finite Representations}
\index{vector!finite representation|textbf}
\index{basis|textbf}
\index{matrix representation|textbf}

\label{ch:training-finite-representations}

\chapterstatus{pedagogical finite-dimensional foundation. The worked examples
are illustrative and synthetic. They establish no material result, convergence
claim, scientific validation, or uncertainty quantification.}

\section{Learning goals}

After working through this chapter, a reader should be able to

\begin{enumerate}
  \item distinguish an abstract vector from a column of its coordinates;
  \item explain span, linear independence, basis, dimension, and
        orthonormality;
  \item distinguish a linear map from the matrix representing it;
  \item read the columns of a matrix as the images of basis vectors;
  \item derive coordinate and matrix transformation rules under a change of
        orthonormal basis;
  \item distinguish coordinate-dependent entries from basis-invariant
        properties; and
  \item state the mathematical and physical metadata required by a finite
        representation.
\end{enumerate}

Chapter~\ref{ch:training-computational-quantities} established that a bare
number is not yet a complete computational quantity. The same lesson now
applies to arrays. A column of numbers becomes the representation of a vector
only after its basis and component meanings are known. A square array becomes
the representation of an operator only after its state space, basis, units,
and conventions are declared.

This chapter assumes only familiar arithmetic with real and complex numbers.
The notation is introduced as it is used. The aim is not abstract algebra for
its own sake. It is to prevent later calculations from silently subtracting,
rotating, diagonalizing, or comparing arrays that represent different physical
objects.

\paragraph{Notation convention.}
An abstract vector is written in italic type, such as \(v\in\mathcal V\). Its
finite coordinate column is bold, such as
\(\mathbf c=[v]_{\mathcal B}\in\mathbb C^N\). Bold symbols such as
\(\mathbf r\), \(\mathbf k\), and \(\mathbf a_i\) are also used for physical
spatial and reciprocal vectors in declared Cartesian coordinates. Matrices use
capital letters such as \(A\), \(H\), and \(U\). Quantum states will later use
ket notation, such as \(\lvert\psi\rangle\). We do not use arrow accents,
because they do not distinguish an abstract state from a finite coordinate
column and become awkward for complex and indexed quantities.

\section{Scalars, vectors, and vector spaces}
\index{scalar}
\index{vector}
\index{vector space}

A \emph{scalar} is an ordinary number used to multiply another mathematical
object. In this monograph scalars are usually complex numbers,
\(\alpha\in\mathbb C\), with real numbers as the special case
\(\alpha\in\mathbb R\).

In elementary mechanics, a vector is often pictured as an arrow with a length
and direction. That picture works well for displacement, velocity, and force,
but it does not cover every vector used in physics. A wavefunction, for
example, is a vector even though it is a function rather than an arrow. The
common mathematical structure is a vector space: its elements can be added to
one another and multiplied by scalars while remaining in the same space.

\begin{definition}[Vector space]
A vector space \(\mathcal V\) is a collection of objects called vectors for
which vector addition and scalar multiplication are defined. If
\(u,v\in\mathcal V\) and \(\alpha,\beta\in\mathbb C\), then
\begin{equation}
  \alpha u+\beta v\in\mathcal V.
\end{equation}
These operations obey the familiar distributive and associative rules of
ordinary coordinate vectors.
\end{definition}

\subsection{Examples of vectors}

The definition is deliberately broad. Its meaning is easier to see through
examples.

\begin{enumerate}
  \item \textbf{A spatial displacement.}
        After Cartesian coordinates are declared, a displacement can be
        represented by
        \begin{equation}
          \mathbf r=\begin{bmatrix}r_x\\r_y\\r_z\end{bmatrix}
          \in\mathbb R^3.
        \end{equation}
        Vector addition combines displacements, and scalar multiplication
        changes their magnitude and possibly their orientation.

  \item \textbf{Orbital amplitudes.}
        In a finite model with \(N\) ordered orbitals, the amplitude column
        \begin{equation}
          \mathbf c
          =\begin{bmatrix}c_1&\cdots&c_N\end{bmatrix}^{\mathsf T}
          \in\mathbb C^N
        \end{equation}
        is a vector. Its entries acquire physical meaning only after the
        orbital ordering and normalization convention are declared.

  \item \textbf{Samples of a function.}
        Values stored at ordered grid points can form
        \begin{equation}
          \mathbf f
          =\begin{bmatrix}f(x_1)&\cdots&f(x_N)\end{bmatrix}^{\mathsf T}.
        \end{equation}
        This finite vector represents sampled information about \(f\); it is
        not the continuous function itself.

  \item \textbf{A polynomial.}
        Polynomials of degree at most \(q\) form a vector space because sums
        and scalar multiples remain polynomials of degree at most \(q\). For
        example,
        \begin{equation}
          p(x)=a_0+a_1x+\cdots+a_qx^q
        \end{equation}
        is an abstract vector in that polynomial space, while
        \(\begin{bmatrix}a_0&\cdots&a_q\end{bmatrix}^{\mathsf T}\) is its
        coordinate column in the ordered basis \((1,x,\ldots,x^q)\).

  \item \textbf{A wavefunction.}
        A wavefunction \(\psi(x)\) is a vector in a function space. It can be
        represented approximately by grid samples or by coefficients in a
        finite basis, but neither finite column is literally identical to the
        continuous function.
\end{enumerate}

The examples illustrate two related uses of the word \emph{vector}. A spatial
displacement, polynomial, or wavefunction is an element of its declared vector
space. A bold finite column is an element of \(\mathbb R^N\) or
\(\mathbb C^N\) and may also serve as the coordinate representation of a
vector in another space. The basis or sampling map supplies that connection.

The symbol \(0_{\mathcal V}\) denotes the zero vector in \(\mathcal V\). It
is the vector satisfying
\begin{equation}
  v+0_{\mathcal V}=v
\end{equation}
for every \(v\in\mathcal V\). Which object plays the role of \(0_{\mathcal V}\)
depends on the space: it may be the zero displacement, the identically zero
function, or a column whose entries all vanish.

\section{Linear combinations, span, and independence}
\index{linear combination}
\index{span}
\index{linear independence}

Given vectors \(v_1,\ldots,v_M\), an expression
\begin{equation}
  c_1v_1+c_2v_2+\cdots+c_Mv_M
\end{equation}
is a \emph{linear combination}. The scalars \(c_j\) are its coefficients.

\begin{definition}[Span]
The span of \(v_1,\ldots,v_M\) is the set of all their linear combinations:
\begin{equation}
  \operatorname{span}\{v_1,\ldots,v_M\}
  =\left\{\sum_{j=1}^{M}c_jv_j\;\middle|\;c_j\in\mathbb C\right\}.
\end{equation}
\end{definition}

For two nonparallel arrows in an ordinary plane, the span is the whole plane.
For two parallel arrows, the span is only one line. Span answers the question,
``Which vectors can these vectors build?''

\begin{definition}[Linear independence]
The vectors \(v_1,\ldots,v_M\) are linearly independent when
\begin{equation}
  \sum_{j=1}^{M}c_jv_j=0_{\mathcal V}
\end{equation}
implies \(c_1=\cdots=c_M=0\). If a nontrivial choice of coefficients gives the
zero vector, the vectors are linearly dependent.
\end{definition}

Linear dependence means that at least one listed vector is redundant. For
example, if \(v_2=2v_1\), then
\begin{equation}
  2v_1-v_2=0_{\mathcal V},
\end{equation}
so \(v_1\) and \(v_2\) are dependent. They describe only one independent
direction even though two vectors were written down.

\section{Bases, dimension, and coordinate columns}
\index{basis!ordered}
\index{dimension}
\index{coordinate vector}

\begin{definition}[Basis]
A basis of a vector space \(\mathcal V\) is a linearly independent set of
vectors that spans \(\mathcal V\).
\end{definition}

Both conditions matter. Spanning ensures that every vector in the space can be
constructed. Independence ensures that the construction is unique.

Suppose
\begin{equation}
  \mathcal B=(b_1,b_2,\ldots,b_N)
\end{equation}
is an ordered basis of a finite-dimensional space \(\mathcal V\). Every
\(v\in\mathcal V\) can be written uniquely as
\begin{equation}
  v=\sum_{i=1}^{N}c_i b_i.
  \label{eq:training-coordinate-expansion}
\end{equation}
The number \(N\) is the \emph{dimension} of the space. The ordered column
\begin{equation}
  [v]_{\mathcal B}
  =\mathbf c
  =\begin{bmatrix}
    c_1\\c_2\\\vdots\\c_N
  \end{bmatrix}
  \in\mathbb C^N
  \label{eq:training-coordinate-column}
\end{equation}
is the coordinate representation of \(v\) in \(\mathcal B\). Here
\(\mathbb C^N\) denotes the set of columns containing \(N\) complex entries.

The brackets \([v]_{\mathcal B}\) mean ``the coordinates of \(v\) in basis
\(\mathcal B\).'' They are useful whenever confusing the vector with its
coordinates would be dangerous.

\subsection{A coordinate example}

Let \(b_1\) point east and \(b_2\) point north, with both having unit length.
The displacement
\begin{equation}
  v=3b_1+b_2
\end{equation}
has coordinate column
\begin{equation}
  [v]_{\mathcal B}
  =\begin{bmatrix}3\\1\end{bmatrix}.
\end{equation}
The column does not intrinsically mean ``three units east and one unit north.''
That meaning comes from the declared basis. In a rotated basis the same
physical displacement has different coordinates.

Ordering is also part of the declaration. If the same basis vectors are listed
as \((b_2,b_1)\), then the coordinate column becomes
\(\begin{bmatrix}1&3\end{bmatrix}^{\mathsf T}\). The physical displacement did
not change; the coordinate convention did.

In software, indices commonly begin at zero even when mathematical labels begin
at one. That implementation convention must not silently change which basis
state a stored component represents.

\section{Inner products, lengths, and orthogonality}
\index{inner product}
\index{norm}
\index{orthogonality}
\index{orthonormal basis}

Vector addition and scalar multiplication tell us how to combine vectors, but
they do not yet tell us how long a vector is or whether two vectors are
perpendicular. In ordinary three-dimensional geometry, that information comes
from the dot product. An inner product generalizes the dot product to complex
coordinate vectors, functions, and the other vector spaces used in physics.

\begin{definition}[Complex inner product]
An inner product \(\langle u,v\rangle\) assigns a complex scalar to two
vectors. We use the convention that it is conjugate-linear in its first
argument and linear in its second. If
\(\mathbf u=[u]_{\mathcal B}\) and
\(\mathbf v=[v]_{\mathcal B}\) are coordinate columns in an orthonormal basis,
then
\begin{equation}
  \langle u,v\rangle
  =\mathbf u^{\dagger}\mathbf v
  =\sum_{i=1}^{N}u_i^*v_i,
  \label{eq:training-coordinate-inner-product}
\end{equation}
where \(^*\) denotes complex conjugation and \(\dagger\) denotes conjugate
transpose.
\end{definition}

The conjugation is essential. It ensures
\begin{equation}
  \langle v,v\rangle
  =\sum_{i=1}^{N}|v_i|^2\geq0,
\end{equation}
with equality only for the zero vector.

The norm induced by the inner product is
\begin{equation}
  \|v\|=\sqrt{\langle v,v\rangle}.
\end{equation}
A vector is \emph{normalized} when \(\|v\|=1\). Two vectors are orthogonal when
\begin{equation}
  \langle u,v\rangle=0.
\end{equation}
A basis is \emph{orthonormal} when
\begin{equation}
  \langle b_i,b_j\rangle=\delta_{ij},
\end{equation}
where \(\delta_{ij}=1\) if \(i=j\) and \(0\) otherwise.

For an orthonormal basis, Equation~\eqref{eq:training-coordinate-expansion} can
be inverted without solving a linear system:
\begin{equation}
  c_i=\langle b_i,v\rangle.
  \label{eq:training-orthonormal-coordinate}
\end{equation}
Taking the inner product with \(b_i\) extracts the component along that basis
vector.

For the remainder of this chapter, all coordinate bases are orthonormal. That
assumption is what permits ordinary coordinate dot products and the simple
coefficient formula above. Nonorthogonal bases require additional information
about basis lengths and angles. Chapter~\ref{ch:training-periodic-geometry}
derives that information through the Gram matrix when it becomes necessary for
nonorthogonal periodic cells.

\section{Linear maps}
\index{linear map}
\index{domain}
\index{codomain}

\begin{definition}[Linear map]
A map \(A:\mathcal V\rightarrow\mathcal W\) is linear when
\begin{equation}
  A(\alpha u+\beta v)
  =\alpha A(u)+\beta A(v)
\end{equation}
for every \(u,v\in\mathcal V\) and all scalars \(\alpha,\beta\).
\end{definition}

The space \(\mathcal V\) is the domain and \(\mathcal W\) is the codomain.
These spaces are part of the definition of the map. Two formulas that produce
the same numbers but act between different declared spaces are not the same
represented map.

Linearity means that knowing the action of \(A\) on a basis determines its
action on every vector. From
\begin{equation}
  v=\sum_{j=1}^{N}c_jb_j
\end{equation}
we obtain
\begin{equation}
  A(v)=\sum_{j=1}^{N}c_jA(b_j).
  \label{eq:training-linearity-on-basis}
\end{equation}
This observation explains why matrix columns have the meaning they do.

\section{Matrices represent linear maps in chosen bases}
\index{matrix!column meaning}
\index{matrix representation}

Let
\begin{equation}
  \mathcal B=(b_1,\ldots,b_N)
\end{equation}
be an ordered orthonormal basis of \(\mathcal V\), and let
\begin{equation}
  \mathcal D=(d_1,\ldots,d_M)
\end{equation}
be an ordered orthonormal basis of \(\mathcal W\). The matrix representing
\(A:\mathcal V\rightarrow\mathcal W\) has entries
\begin{equation}
  [A]_{\mathcal D\leftarrow\mathcal B,ij}
  =\langle d_i,A(b_j)\rangle.
  \label{eq:training-map-matrix-elements}
\end{equation}
It has \(M\) rows and \(N\) columns.

The arrow in the subscript records the coordinate direction: input coordinates
are taken in \(\mathcal B\), and output coordinates are returned in
\(\mathcal D\). Column \(j\) contains all coordinates of \(A(b_j)\):
\begin{equation}
  A(b_j)
  =\sum_{i=1}^{M}
   [A]_{\mathcal D\leftarrow\mathcal B,ij}d_i.
\end{equation}
Therefore, if \(\mathbf x=[v]_{\mathcal B}\), then
\begin{equation}
  [A(v)]_{\mathcal D}
  =[A]_{\mathcal D\leftarrow\mathcal B}\mathbf x.
  \label{eq:training-matrix-action}
\end{equation}

\subsection{Reading a matrix by columns}

Suppose a map on a two-dimensional space satisfies
\begin{equation}
  A(b_1)=2b_1+b_2,
  \qquad
  A(b_2)=-b_1+3b_2.
\end{equation}
Its matrix in \(\mathcal B=(b_1,b_2)\) is
\begin{equation}
  [A]_{\mathcal B}
  =\begin{bmatrix}
    2 & -1\\
    1 & 3
  \end{bmatrix}.
\end{equation}
The first column is the coordinate column of \(A(b_1)\); the second is the
coordinate column of \(A(b_2)\).

For
\begin{equation}
  v=4b_1+2b_2,
  \qquad
  [v]_{\mathcal B}=\begin{bmatrix}4\\2\end{bmatrix},
\end{equation}
matrix multiplication gives
\begin{equation}
  [A(v)]_{\mathcal B}
  =\begin{bmatrix}
    2 & -1\\
    1 & 3
  \end{bmatrix}
  \begin{bmatrix}4\\2\end{bmatrix}
  =\begin{bmatrix}6\\10\end{bmatrix}.
\end{equation}
The multiplication combines the two matrix columns with the same coefficients
used to build \(v\), exactly as Equation~\eqref{eq:training-linearity-on-basis}
requires.

A square matrix represents a map from a finite space to itself. A rectangular
matrix represents a map between spaces of different dimensions or between
different selected subspaces. Shape alone does not identify what either axis
means.

\section{Identity, adjoint, and unitary maps}
\index{identity map}
\index{adjoint}
\index{unitary matrix}

The identity map leaves every vector unchanged:
\begin{equation}
  I(v)=v.
\end{equation}
In any basis its matrix is the identity matrix \(I_N\), with ones on the
diagonal and zeros elsewhere.

For a matrix \(A\), the conjugate transpose \(A^{\dagger}\) is obtained by
transposing the array and complex-conjugating every entry. It represents the
adjoint map when the domain and codomain bases are orthonormal. The adjoint is
characterized abstractly by
\begin{equation}
  \langle w,A v\rangle
  =\langle A^{\dagger}w,v\rangle.
\end{equation}
For real matrices, conjugate transpose reduces to ordinary transpose.

\begin{definition}[Unitary matrix]
A square complex matrix \(U\) is unitary when
\begin{equation}
  U^{\dagger}U=UU^{\dagger}=I_N.
  \label{eq:training-unitary-definition}
\end{equation}
Equivalently, \(U^{-1}=U^{\dagger}\).
\end{definition}

The columns of a unitary matrix are an orthonormal basis of \(\mathbb C^N\).
Unitary matrices preserve coordinate inner products and Euclidean norms:
\begin{equation}
  (U\mathbf u)^{\dagger}(U\mathbf v)
  =\mathbf u^{\dagger}\mathbf v,
  \qquad
  \|U\mathbf v\|_2=\|\mathbf v\|_2.
\end{equation}
Real unitary matrices are called orthogonal matrices.

Chapter~\ref{ch:training-hermitian-eigensystems} will define Hermitian matrices
and show why unitary basis matrices are central to their diagonalization.

\section{Changing coordinates without changing the object}
\index{basis transformation}
\index{coordinate transformation}

Let
\begin{equation}
  \mathcal B=(b_1,\ldots,b_N)
\end{equation}
and
\begin{equation}
  \mathcal B'=(b'_1,\ldots,b'_N)
\end{equation}
be two ordered orthonormal bases of the same space. Store the new basis vectors
as columns in the old basis:
\begin{equation}
  b'_j=\sum_{i=1}^{N}b_iU_{ij}.
  \label{eq:training-new-basis-columns}
\end{equation}
The matrix \(U\) is unitary because both bases are orthonormal.

Let \(\mathbf c=[v]_{\mathcal B}\) and
\(\mathbf c'=[v]_{\mathcal B'}\). Expanding the same abstract vector in both
bases gives
\begin{equation}
  v=\sum_i c_i b_i=\sum_j c'_j b'_j.
\end{equation}
Substitute Equation~\eqref{eq:training-new-basis-columns} into the second sum.
The old coordinate column then satisfies
\begin{equation}
  \mathbf c=U\mathbf c'.
\end{equation}
Multiplication by \(U^{\dagger}=U^{-1}\) gives
\begin{equation}
  \mathbf c'=U^{\dagger}\mathbf c.
  \label{eq:training-vector-coordinate-change}
\end{equation}
The vector did not move. Only its coordinate description changed. This is a
\emph{passive} change of basis. It should not be confused with actively applying
a physical rotation or symmetry operation to the vector.

If \(A:\mathcal V\rightarrow\mathcal V\), then its matrix in the new basis is
\begin{equation}
  [A]_{\mathcal B'}
  =U^{\dagger}[A]_{\mathcal B}U.
  \label{eq:training-operator-basis-change}
\end{equation}
The factor on the right converts new input coordinates to old coordinates. The
factor on the left converts the old output coordinates back to new coordinates.

For a map between two different spaces, the domain and codomain can have
different basis changes. If \(U_{\mathcal V}\) changes the domain basis and
\(U_{\mathcal W}\) changes the codomain basis, then
\begin{equation}
  [A]'
  =U_{\mathcal W}^{\dagger}[A]U_{\mathcal V}.
  \label{eq:training-rectangular-basis-change}
\end{equation}
This more general rule matters whenever a rectangular projection, embedding,
or comparison map is represented.

\subsection{A rotated-coordinate example}

Define
\begin{equation}
  b'_1=\frac{b_1+b_2}{\sqrt2},
  \qquad
  b'_2=\frac{-b_1+b_2}{\sqrt2}.
\end{equation}
Then
\begin{equation}
  U=\frac{1}{\sqrt2}
  \begin{bmatrix}
    1 & -1\\
    1 & 1
  \end{bmatrix}.
\end{equation}
For the earlier vector \(v=3b_1+b_2\),
\begin{equation}
  [v]_{\mathcal B'}
  =U^{\mathsf T}[v]_{\mathcal B}
  =\frac{1}{\sqrt2}
   \begin{bmatrix}
     1 & 1\\
     -1 & 1
   \end{bmatrix}
   \begin{bmatrix}3\\1\end{bmatrix}
  =\begin{bmatrix}2\sqrt2\\-\sqrt2\end{bmatrix}.
\end{equation}
The columns \(\begin{bmatrix}3&1\end{bmatrix}^{\mathsf T}\) and
\(\begin{bmatrix}2\sqrt2&-\sqrt2\end{bmatrix}^{\mathsf T}\) look different,
but they represent the same vector in different bases.

\section{What changes and what does not}
\index{basis invariant}
\index{coordinate dependent}

Under the unitary similarity transformation in
Equation~\eqref{eq:training-operator-basis-change}, individual matrix entries
generally change. The following quantities do not:

\begin{itemize}
  \item dimension and rank;
  \item trace and determinant;
  \item eigenvalues, including multiplicities; and
  \item unitarily invariant norms, including the Frobenius norm.
\end{itemize}

These are basis invariants. In contrast, a particular diagonal entry, a named
row, or a componentwise residual depends on the basis unless additional
structure identifies those coordinates.

A diagonal matrix is not intrinsically diagonal in every basis. It states that
the selected coordinate axes are adapted to the map. After a general basis
change, the same map may have off-diagonal entries. Conversely, making a matrix
look different by changing coordinates does not by itself create a different
physical operator.

Basis invariance also has limits. Two arrays with equal eigenvalues are not
automatically representations of the same physical operator. Their state
spaces, basis meanings, energy references, and other physical metadata must
still be identified. Chapter~\ref{ch:training-hermitian-eigensystems} develops
this distinction for spectra and eigenspaces.

\section{Finite representations and their contracts}
\index{finite representation!contract}

\begin{definition}[Finite representation]
A finite representation is a finite array together with the mathematical and
physical declarations needed to interpret its indices and permitted
operations.
\end{definition}

A useful representation contract answers at least four questions.

\begin{enumerate}
  \item \textbf{What is represented?} Name the vector, linear map, bilinear
        form, or operator and its source model.
  \item \textbf{On which finite spaces?} State the domain, codomain,
        dimensions, and any selected subspaces.
  \item \textbf{In which coordinates?} State ordered bases, axis meanings,
        index ordering, normalization, and gauge conventions.
  \item \textbf{With which quantities?} State units, physical references,
        geometry, and provenance required by the promised operations.
\end{enumerate}

A finite matrix can arise in several inequivalent ways:

\begin{itemize}
  \item by declaring a model with finitely many degrees of freedom;
  \item by restricting an existing map to a finite invariant subspace;
  \item by projecting onto a selected subspace;
  \item by sampling a continuous kernel or function; or
  \item by discretizing a differential operator.
\end{itemize}

These constructions are not synonyms. A \(100\times100\) matrix does not reveal
whether it is an exact finite model, a projection, or a grid approximation.
That distinction belongs to the representation contract.

Sparse and dense arrays may represent the same finite map, but their
computational contracts differ. Sparse structure should be preserved whenever
the mathematics and downstream operations allow it. Converting a sparse
operator to a dense array merely to test symmetry or Hermiticity can turn a
local diagnostic into an unnecessary memory-intensive operation.

\section{A two-basis calculation}

We first perform the complete calculation symbolically. Let a vector and a
linear map have the following representations in the original basis
\(\mathcal B\):
\begin{equation}
  \mathbf c=[v]_{\mathcal B}
  =\begin{bmatrix}3\\1\end{bmatrix},
  \qquad
  A_{\mathcal B}
  =\begin{bmatrix}
    2 & -1\\
    1 & 3
  \end{bmatrix}.
  \label{eq:training-two-basis-original-data}
\end{equation}
Use the rotated basis whose columns in the original coordinates form
\begin{equation}
  U=\frac{1}{\sqrt2}
  \begin{bmatrix}
    1 & -1\\
    1 & 1
  \end{bmatrix},
  \qquad
  U^{\mathsf T}U=I_2.
  \label{eq:training-two-basis-rotation}
\end{equation}
Because this example is real, \(U^{\dagger}=U^{\mathsf T}\). The coordinates of
the same vector in the new basis are
\begin{equation}
  \mathbf c'
  =U^{\mathsf T}\mathbf c
  =\begin{bmatrix}
    2\sqrt2\\
    -\sqrt2
  \end{bmatrix}.
  \label{eq:training-two-basis-vector-result}
\end{equation}
The matrix representing the same linear map in the new basis is
\begin{align}
  A_{\mathcal B'}
  &=U^{\mathsf T}A_{\mathcal B}U\\
  &=\begin{bmatrix}
      5/2 & -1/2\\
      3/2 & 5/2
    \end{bmatrix}.
  \label{eq:training-two-basis-matrix-result}
\end{align}
Neither \(\mathbf c'=\mathbf c\) nor
\(A_{\mathcal B'}=A_{\mathcal B}\). That is expected:
the coordinate descriptions changed.

We can confirm that the represented action did not change. In the original
basis,
\begin{equation}
  [A(v)]_{\mathcal B}
  =A_{\mathcal B}\mathbf c
  =\begin{bmatrix}5\\6\end{bmatrix}.
\end{equation}
In the new basis,
\begin{align}
  [A(v)]_{\mathcal B'}
  &=A_{\mathcal B'}\mathbf c'\\
  &=\begin{bmatrix}
      11/\sqrt2\\
      1/\sqrt2
    \end{bmatrix}\\
  &=U^{\mathsf T}[A(v)]_{\mathcal B}.
  \label{eq:training-two-basis-action-result}
\end{align}
The two output columns differ because they use different bases, but the final
equality shows that they are coordinates of the same output vector.

The code below repeats these steps numerically. It prints the transformed
objects and small round-trip errors so that a reader can inspect what happened.
It does not use a testing framework.

\begin{pythonexample}{representing one linear map in two bases}
import numpy as np

# Coordinates of one abstract vector in the original ordered basis.
c = np.array([3.0, 1.0])

# Columns of U are the new orthonormal basis vectors in old coordinates.
U = np.array([[1.0, -1.0], [1.0, 1.0]]) / np.sqrt(2.0)

# Represent one synthetic linear map in the original basis.
A = np.array([[2.0, -1.0], [1.0, 3.0]])

# Describe the same vector and map in the new basis.
c_prime = U.T @ c
A_prime = U.T @ A @ U

# Compute the map's action in both coordinate systems.
Av_old = A @ c
Av_new = A_prime @ c_prime

print("coordinates in the old basis:", c)
print("coordinates in the new basis:", c_prime)
print("matrix in the old basis:\n", A)
print("matrix in the new basis:\n", A_prime)
print("output coordinates in the old basis:", Av_old)
print("output coordinates in the new basis:", Av_new)

# Values near zero show agreement up to floating-point roundoff.
print("basis orthonormality error:", np.linalg.norm(U.T @ U - np.eye(2)))
print("vector round-trip error:", np.linalg.norm(U @ c_prime - c))
print("matrix round-trip error:", np.linalg.norm(U @ A_prime @ U.T - A))
print("action consistency error:", np.linalg.norm(Av_new - U.T @ Av_old))
\end{pythonexample}

This example demonstrates coordinate transformation rules on one synthetic
vector and one synthetic \(2\times2\) matrix. The displayed errors are
floating-point consistency diagnostics, not physical uncertainties. The
example does not identify a physical state space, operator, Hamiltonian, or
material system.

\documentcallout{Notebook to do}{Create
\path{examples/tutorials/research-monograph/foundations/finite_representations.ipynb}
with explicit imports, explanatory inline comments, span and independence
examples, ordered-basis metadata, rectangular maps, orthonormal basis changes,
and sparse-preserving invariant checks. Keep it unexecuted in the repository.}

\section{Common mistakes}

\begin{itemize}
  \item Treating an abstract vector as though it were identical to one of its
        coordinate columns.
  \item Calling a set a basis after checking only that it spans, or only that
        it is independent.
  \item Forgetting that basis ordering determines coordinate ordering.
  \item Reading matrix rows as the images of input basis vectors; those images
        are stored in columns under the convention used here.
  \item Applying \(U\) where a passive coordinate change requires
        \(U^{\dagger}\).
  \item Comparing entries of two matrices before identifying their domain,
        codomain, and bases.
  \item Assuming equal shape or equal eigenvalues establishes equal represented
        meaning.
  \item Densifying a sparse representation merely to run a diagnostic.
\end{itemize}

\section{From representations to eigensystems}

A matrix may admit a basis in which its action becomes diagonal. Finding and
interpreting that basis is not merely another array conversion: existence,
orthogonality, multiplicity, and numerical stability depend on the class of
matrix being studied. Chapter~\ref{ch:training-hermitian-eigensystems} develops
those questions for the real-symmetric and complex-Hermitian matrices used
throughout the monograph.

\section{Exercises}

\begin{enumerate}
  \item Decide whether each of the following is a vector, a coordinate
        representation, or either depending on context: a velocity arrow, a
        list of grid values, a polynomial, and a NumPy array.
  \item Determine whether
        \(\{(1,0),(2,0)\}\) spans \(\mathbb R^2\), and determine whether it is
        linearly independent.
  \item Show that the coordinate coefficients in an orthonormal basis satisfy
        \(c_i=\langle b_i,v\rangle\).
  \item Derive Equation~\eqref{eq:training-vector-coordinate-change} from
        Equation~\eqref{eq:training-new-basis-columns}.
  \item Derive Equation~\eqref{eq:training-operator-basis-change} by following
        an input column from new coordinates to old coordinates, applying the
        old matrix, and converting the result back.
  \item Give an example of a rectangular matrix and identify its domain basis,
        codomain basis, row meaning, and column meaning.
  \item Construct two matrices representing the same map in different
        orthonormal bases. Identify entries that change and quantities that
        remain invariant.
  \item Give one example each of a finite matrix produced by restriction,
        projection, sampling, and discretization.
\end{enumerate}

\section{Evidence boundary}

Exact coordinate round trips and invariant checks can verify implementation of
a finite representation. They do not establish that the represented space is
physically appropriate, that a discretization has converged, that metadata are
truthful, or that the matrix models a real material. Those are separate claims
requiring model, provenance, numerical, and validation evidence.
